\section{Stable rank-tight occupation}\label{sec:stable51}
The previous volume budgets concern exact machines. We now derive a
uniform perturbation estimate from the prescribed signed seeds and
bounded decoders. No lower bound on individual hidden-state
probabilities is assumed. The structural hypotheses remain those of
Section~\ref{sec:lifts51}. Write $\delta=2\varepsilon$ for mean
error corresponding to binary total-variation error $\varepsilon$.

\begin{lemma}[Quantitative recovery from anchor distributions]
\label{lem:anchor51}
Suppose a machine has worst-word mean error at most $\delta$, where
$0\le\delta<\rho/(2D)$. Every command cut has width at least $D+1$.
At a cut of width exactly $D+1$, let $Z_t$ be its actual identity-suffix
mean matrix and let $B_t(v)$ be its actual mean matrix for a suffix
$v$. Then
\begin{equation}\label{eq:row-stability51}
 \max_{k,j}|B_t(v)_{kj}-(Z_tU_v^T)_{kj}|
 \le h,\qquad
 h=\frac{2(D+1)(1+\sqrt D)\delta}{\rho-2D\delta}.
\end{equation}
This bound holds for every suffix, independently of its length and of
all widths at later cuts. Moreover $S_t=\conv\{Z_t(k,\cdot)^T\}$
is a full-dimensional simplex and
\begin{equation}\label{eq:inner-stability51}
 C_b\subset S_t\subset[-1,1]^D,\qquad b=\rho-D\delta.
\end{equation}
\end{lemma}
\begin{proof}
Let $p_{i,+}$ and $p_{i,-}$ be the actual state distributions after
seed $\pm e_i$ and an identity prefix of length $t$. Form a
$(D+1)$-by-$K_t$ stochastic matrix $F_t$ whose first row is
$(p_{1,+}+p_{1,-})/2$ and whose other rows are $p_{i,+}$,
$1\le i\le D$. Put $A_t=[\mathbf1,Z_t]$ and $M_t=F_tA_t$.
The target matrix and its inverse are
\[
 M_* =\begin{pmatrix}1&0\\\mathbf1&\rho I_D\end{pmatrix},
 \qquad
 M_*^{-1}=\begin{pmatrix}1&0\\-\rho^{-1}\mathbf1&\rho^{-1}I_D\end{pmatrix}.
\]
Here $\|M\|_\infty=\max_i\sum_j|M_{ij}|$ is the induced infinity
norm, distinct from the entrywise maximum in
\eqref{eq:row-stability51}. The averaged first anchor has error at most
$\delta$, since the average of two errors of size at most $\delta$
still has size at most $\delta$. In the induced infinity norm,
$\|M_*^{-1}\|_\infty=2/\rho$ and
$\|M_t-M_*\|_\infty\le D\delta$. The first column is exact;
each other entry has error at most $\delta$. The Neumann series
therefore proves invertibility and
\[
 \|M_t^{-1}\|_\infty\le\frac{2}{\rho-2D\delta}.
\]
Since $M_t$ has rank $D+1$, necessarily $K_t\ge D+1$.
When $K_t=D+1$, both $F_t$ and $A_t$ are invertible and
\begin{equation}\label{eq:anchor-inverse51}
 \|F_t^{-1}\|_\infty
 =\|A_tM_t^{-1}\|_\infty
 \le\frac{2(D+1)}{\rho-2D\delta}.
\end{equation}

For the same anchor prefixes, the target means after suffix $v$ are
the identity-suffix target means multiplied by $U_v^T$. The actual
errors of each of these two sets of means are at most $\delta$ per
coordinate. Each row of an orthogonal matrix has absolute sum at most
$\sqrt D$, so every entry of
$F_t(B_t(v)-Z_tU_v^T)$ has absolute value at most
$(1+\sqrt D)\delta$. Multiply on the left by $F_t^{-1}$ and use
\eqref{eq:anchor-inverse51}. This proves
\eqref{eq:row-stability51}; no transition-by-transition error sum has
been used.

The points $p_{i,\pm}Z_t$ lie in $S_t$ and are within $\delta$ in
each coordinate of $\pm\rho e_i^T$. Consequently the support
function of $S_t$ satisfies, for every $u\in\R^D$,
\[
 h_{S_t}(u)\ge\rho\|u\|_\infty-\delta\|u\|_1
             \ge(\rho-D\delta)\|u\|_\infty=h_{C_b}(u).
\]
Support-function separation gives $C_b\subset S_t$. The outer
bound follows from $|Z_t(k,j)|\le1$. Since $b>0$, $S_t$ has full
dimension; its $D+1$ generators are therefore affinely independent.
\end{proof}

\begin{theorem}[Robust occupation with separated narrow cuts]
\label{thm:robust51}
Assume $I,-I\in\calA$ and $0\le\delta<\rho/(2D)$. Define
\begin{equation}\label{eq:Lambda51}
 b=\rho-D\delta,\qquad
 \Lambda=1+\frac{2D(D+1)(1+\sqrt D)\delta}
                    {(\rho-2D\delta)(\rho-D\delta)}.
\end{equation}
For a machine of mean error at most $\delta$, let $m\ge1$ be the
number of command cuts of width $D+1$, including either endpoint
when it has this width. Then
\begin{equation}\label{eq:robust-budget51}
 \left(\frac{c_D}{\Lambda^D}\right)^{m-1}
       \le D!b^{-D},\qquad c_D=2^{-D}{2D\choose D}.
\end{equation}
All other command cuts may have arbitrary finite widths. For every
pair $s<t$ of selected cuts and every word $w$ of length $t-s$,
\begin{equation}\label{eq:approx-containment51}
 U_wS_s\subset S_t+h[-1,1]^D\subset\Lambda S_t,
\end{equation}
where $h$ is as in \eqref{eq:row-stability51}.
\end{theorem}
\begin{proof}
Use the actual composed stochastic matrix $T_w$ between the two cuts.
Its product with $Z_t$ is the conditional mean matrix at $s$ for the
suffix consisting of $w$ followed by identities. Each row of $T_wZ_t$
is a convex combination of the rows of $Z_t$.
Lemma~\ref{lem:anchor51} at $s$ therefore puts the image of every
vertex of $S_s$ under $U_w$ within $h[-1,1]^D$ of $S_t$.
Convexity proves the first containment.

Since $C_b\subset S_t$, the cube $h[-1,1]^D$ lies in
$(Dh/b)C_b\subset(Dh/b)S_t$. Minkowski addition and convexity give
$S_t+(Dh/b)S_t=(1+Dh/b)S_t=\Lambda S_t$, proving the second
containment. Between consecutive selected cuts both the identity word
and a word containing one $-I$ and otherwise identities are available.
Thus $\Lambda S_t$ contains $S_s$ and $-S_s$.
Lemma~\ref{lem:difference50} gives
\[
 \operatorname{vol}(S_t)\ge c_D\Lambda^{-D}
                                  \operatorname{vol}(S_s).
\]
Iterate over the $m-1$ intervals and use
$2^Db^D/D!\le\operatorname{vol}(S_{\rm first})$ and
$\operatorname{vol}(S_{\rm last})\le2^D$.
This proves \eqref{eq:robust-budget51}.
\end{proof}

At $\delta=0$ the exact budget is recovered. For fixed $D,\rho$,
the multiplier is greater than one for all sufficiently small
positive $\delta$, with its range given explicitly by
\eqref{eq:Lambda51}. Its dependence on $\rho$ is essential: a
vanishing signal cannot supply a uniform affine-conditioning bound.
The next consequence gives simple rational constants in dimension two.

\begin{corollary}[A horizon-uniform positive-error frontier]
\label{cor:robust-four51}
Let $D=2$ and let $\calA$ be any signed-permutation alphabet containing
$I,-I$. For every $0<\rho<1$,
\begin{equation}\label{eq:robust-four51}
 0\le\varepsilon\le\frac{\rho^2}{2000},\qquad
 \left(\frac{512}{363}\right)^N>
       2\left(\frac{500}{499}\right)^2\rho^{-2}
 \quad\Longrightarrow\quad W_{N,\varepsilon}=4.
\end{equation}
In particular, when $\rho=1/10$, all $N\ge16$ and all
$0\le\varepsilon\le1/200000$ satisfy the conclusion. These are
sufficient constants, not the sharp error threshold or the first
four-state horizon.
\end{corollary}
\begin{proof}
Put $\delta=2\varepsilon\le\rho^2/1000$. Then
$b=\rho-2\delta\ge(499/500)\rho$. Since $\rho<1$ and
$\sqrt2<3/2$, equation~\eqref{eq:Lambda51} gives
\[
 \Lambda-1\le
 \frac{30/1000}{(996/1000)(998/1000)}<\frac1{32};
\]
the last comparison is $960000<994008$. Hence
$c_2/\Lambda^2\ge(3/2)(32/33)^2=512/363$.
A peak of at most three would, by Lemma~\ref{lem:anchor51}, have
exactly three labels at every command cut. Apply
\eqref{eq:robust-budget51} with $m=N+1$ to contradict the strict
inequality in \eqref{eq:robust-four51}. Four labels attain error
zero by the signed-coordinate permutation construction of
Corollary~\ref{cor:four50}.
For the numerical claim, the rational inequality
$(512/363)^{16}>200(500/499)^2$ verifies the horizon bound;
monotonicity gives all larger horizons. The error bound is
$(1/10)^2/2000=1/200000$.
\end{proof}

The upper error bound in this interval does not shrink as the horizon
grows; the sufficient horizon is a separate condition. It strengthens,
for this class, the fixed-horizon compactness threshold
$\eta_N(\calA,\rho)$ in Corollary~\ref{cor:four50}. This theorem controls approximate rank-tight cuts. The unrestricted
positive-error finite-group classification is proved separately in
Theorem~\ref{thm:noisy-finitegroup52}; its proof uses finite-orbit
reduction rather than extending the anchor inverse to nonsquare matrices.
